\textbf{Phần 1. Trắc nghiệm lựa chọn}
\setcounterpageref{ex}{0}

\begin{ex} %[2H3N18]
Đơn vị dùng để đo khối lượng là
\choice
{\True Ki-lô-gam}
{Lít}
{Mét}
{Đề-xi-mét}
\loigiai{
Đơn vị đo khối lượng là ki-lô-gam.\\
\textbf{Đáp án của câu hỏi là A}
}
\end{ex}

\begin{ex} %[2H3N18]
Đơn vị dùng để đo sức chứa chất lỏng là
\choice
{Ki-lô-gam}
{Xăng-ti-mét}
{Giờ}
{\True Lít}
\loigiai{
Đo dung tích dùng đơn vị lít.\\
\textbf{Đáp án của câu hỏi là D}
}
\end{ex}

\begin{ex} %[2H3H18]
Tính: $25\text{ kg} + 15\text{ kg} = \dots$
\choice
{$30\text{ kg}$}
{$50\text{ kg}$}
{\True $40\text{ kg}$}
{$35\text{ kg}$}
\loigiai{
$25 + 15 = 40\text{ kg}$.\\
\textbf{Đáp án của câu hỏi là C}
}
\end{ex}

\begin{ex} %[2H3H18]
Tính: $48\text{ l} - 20\text{ l} = \dots$
\choice
{$26\text{ l}$}
{\True $28\text{ l}$}
{$30\text{ l}$}
{$18\text{ l}$}
\loigiai{
$48 - 20 = 28\text{ l}$.\\
\textbf{Đáp án của câu hỏi là B}
}
\end{ex}

\begin{ex} %[2H3H18]
Bao gạo nặng $45\text{ kg}$, bao ngô nhẹ hơn bao gạo $10\text{ kg}$. Bao ngô nặng số ki-lô-gam là
\choice
{$30\text{ kg}$}
{$40\text{ kg}$}
{$55\text{ kg}$}
{\True $35\text{ kg}$}
\loigiai{
$45 - 10 = 35\text{ kg}$.\\
\textbf{Đáp án của câu hỏi là D}
}
\end{ex}

\begin{ex} %[2H3H18]
Can to có $15\text{ l}$ dầu, can nhỏ có $8\text{ l}$ dầu. Cả hai can có số lít dầu là
\choice
{\True $23\text{ l}$}
{$22\text{ l}$}
{$24\text{ l}$}
{$25\text{ l}$}
\loigiai{
$15 + 8 = 23\text{ l}$.\\
\textbf{Đáp án của câu hỏi là A}
}
\end{ex}

\begin{ex} %[2H3V18]
Tính giá trị biểu thức: $50\text{ kg} - 20\text{ kg} + 15\text{ kg}$.
\choice
{$40\text{ kg}$}
{\True $45\text{ kg}$}
{$50\text{ kg}$}
{$55\text{ kg}$}
\loigiai{
$30 + 15 = 45\text{ kg}$.\\
\textbf{Đáp án của câu hỏi là B}
}
\end{ex}

\begin{ex} %[2H3V18]
Tính giá trị biểu thức: $60\text{ l} - 15\text{ l} - 15\text{ l}$.
\choice
{\True $30\text{ l}$}
{$25\text{ l}$}
{$35\text{ l}$}
{$40\text{ l}$}
\loigiai{
$60 - 30 = 30\text{ l}$.\\
\textbf{Đáp án của câu hỏi là A}
}
\end{ex}

\begin{ex} %[2H3V18]
Mẹ mua $12\text{ kg}$ dưa hấu và $8\text{ kg}$ xoài. Mẹ mua tất cả số ki-lô-gam quả là
\choice
{$18\text{ kg}$}
{$22\text{ kg}$}
{\True $20\text{ kg}$}
{$24\text{ kg}$}
\loigiai{
$12 + 8 = 20\text{ kg}$.\\
\textbf{Đáp án của câu hỏi là C}
}
\end{ex}

\begin{ex} %[2H3V18]
Bình thứ nhất có $14\text{ l}$ nước, bình thứ hai có ít hơn bình thứ nhất $5\text{ l}$ nước. Bình thứ hai có số lít nước là
\choice
{$8\text{ l}$}
{$10\text{ l}$}
{$11\text{ l}$}
{\True $9\text{ l}$}
\loigiai{
$14 - 5 = 9\text{ l}$.\\
\textbf{Đáp án của câu hỏi là D}
}
\end{ex}

\begin{ex} %[2H3V18]
Điền dấu thích hợp: $35\text{ kg} + 10\text{ kg} \dots 50\text{ kg}$.
\choice
{$>$}
{\True $<$}
{$=$}
{Không so sánh được}
\loigiai{
$35 + 10 = 45\text{ kg} < 50\text{ kg}$.\\
\textbf{Đáp án của câu hỏi là B}
}
\end{ex}

\begin{ex} %[2H3C18]
Có hai bao đường, bao thứ nhất nặng $25\text{ kg}$. Nếu chuyển $5\text{ kg}$ từ bao thứ nhất sang bao thứ hai thì hai bao nặng bằng nhau. Hỏi bao thứ hai ban đầu nặng bao nhiêu ki-lô-gam?
\choice
{$10\text{ kg}$}
{$20\text{ kg}$}
{\True $15\text{ kg}$}
{$25\text{ kg}$}
\loigiai{
Sau khi chuyển, bao 1 nặng: $25 - 5 = 20\text{ kg}$.\\
Lúc đó bao 2 cũng nặng $20\text{ kg} \Rightarrow$ Ban đầu bao 2 nặng: $20 - 5 = 15\text{ kg}$.\\
\textbf{Đáp án của câu hỏi là C}
}
\end{ex}

\textbf{Phần 2. Trắc nghiệm đúng sai}
\setcounterpageref{ex}{0}

\begin{ex} %[2H3H18]
Xét tính đúng sai về khối lượng và dung tích:
\choiceTF
{\True Ki-lô-gam viết tắt là $\text{kg}$.}
{\True Lít viết tắt là $\text{l}$.}
{$20\text{ kg} + 10\text{ l} = 30\text{ kg}$.}
{\True Không thể cộng hai đại lượng khác đơn vị đo ($	ext{kg}$ và $	ext{l}$).}
\loigiai{
\begin{itemize}
    \item $\text{kg}$ là ki-lô-gam: Đúng.
    \item $\text{l}$ là lít: Đúng.
    \item Khác đơn vị đo không cộng trực tiếp được. Ý 3 Sai, Ý 4 Đúng.
\end{itemize}
\begin{itemize}
    \item a) \textbf{Đúng}.
    \item b) \textbf{Đúng}.
    \item c) \textbf{Sai}.
    \item d) \textbf{Đúng}.
\end{itemize}
}
\end{ex}

\begin{ex} %[2H3V18]
Thùng A có $35\text{ l}$ dầu, thùng B có $25\text{ l}$ dầu. Xét tính đúng sai:
\choiceTF
{\True Cả hai thùng có tất cả $60\text{ l}$ dầu.}
{\True Thùng A nhiều hơn thùng B là $10\text{ l}$ dầu.}
{Nếu múc $10\text{ l}$ từ thùng A thì hai thùng bằng nhau.}
{\True Khi đó thùng A còn $35 - 10 = 25\text{ l}$ bằng thùng B.}
\loigiai{
\begin{itemize}
    \item Tổng: $35 + 25 = 60\text{ l}$: Đúng.
    \item Hiệu: $35 - 25 = 10\text{ l}$: Đúng.
    \item Múc 10 lít thì còn 25 lít: Đúng. Cả hai ý cuối đều Đúng.
\end{itemize}
\begin{itemize}
    \item a) \textbf{Đúng}.
    \item b) \textbf{Đúng}.
    \item c) \textbf{Đúng}.
    \item d) \textbf{Đúng}.
\end{itemize}
}
\end{ex}

\textbf{Phần 3. Trả lời ngắn}
\setcounterpageref{ex}{0}

\begin{ex} %[2H3H18]
Tính: $28\text{ kg} + 12\text{ kg}$.
\shortans{40}
\loigiai{
\textbf{Đáp số:} $40$\\
$28 + 12 = 40\text{ kg}$.
}
\end{ex}

\begin{ex} %[2H3H18]
Tính: $50\text{ l} - 18\text{ l}$.
\shortans{32}
\loigiai{
\textbf{Đáp số:} $32$\\
$50 - 18 = 32\text{ l}$.
}
\end{ex}

\begin{ex} %[2H3V18]
Con lợn to nặng $75\text{ kg}$, con lợn nhỏ nặng $45\text{ kg}$. Con lợn to nặng hơn con lợn nhỏ bao nhiêu ki-lô-gam?
\shortans{30}
\loigiai{
\textbf{Đáp số:} $30$\\
$75 - 45 = 30\text{ kg}$.
}
\end{ex}

\begin{ex} %[2H3V18]
Một can chứa $20\text{ l}$ dầu, mẹ đã dùng $12\text{ l}$ dầu. Hỏi trong can còn lại bao nhiêu lít dầu?
\shortans{8}
\loigiai{
\textbf{Đáp số:} $8$\\
$20 - 12 = 8\text{ l}$.
}
\end{ex}

\begin{ex} %[2H3V18]
Tính giá trị biểu thức: $25\text{ kg} + 35\text{ kg} - 20\text{ kg}$.
\shortans{40}
\loigiai{
\textbf{Đáp số:} $40$\\
$60 - 20 = 40\text{ kg}$.
}
\end{ex}

\begin{ex} %[2H3C18]
Thùng thứ nhất có $24\text{ l}$ sữa, thùng thứ hai có nhiều hơn thùng thứ nhất $6\text{ l}$ sữa. Hỏi cả hai thùng có bao nhiêu lít sữa?
\shortans{54}
\loigiai{
\textbf{Đáp số:} $54$\\
Thùng 2: $24 + 6 = 30\text{ l}$. Cả hai thùng: $24 + 30 = 54\text{ l}$.
}
\end{ex}
